\documentclass[a4,11pt]{aleph-notas}

% -- Paquetes adicionales
\usepackage{enumitem}
\usepackage{aleph-comandos}

% -- Datos
\institucion{Facultad de Ciencias Exactas, Naturales y Ambientales}
\carrera{Catálogo STEM}
\asignatura{Matemática III}
\tema{Resumen 13: Integrales triples}
\autor[A. Merino]{Andrés Merino}
\fecha{Periodo 2026-1}

\logouno[0.16\textwidth]{Logos/logoPUCE_04_ac}
\definecolor{colortext}{HTML}{1957d1}
\definecolor{colordef}{HTML}{1957d1}
\fuente{montserrat}


\begin{document}

\encabezado

En esta sección asumiremos $a,b,c,d,p,q\in\R$ tal que $a<b$, $c<d$ y $p<q$ y definimos
\[
    B=[a,b]\times [c,d] \times [p,q]
\]
al cual llamaremos un bloque.

%%%%%%%%%%%%%%%%%%%%%%%%%%%%%%%%%%%%%%
\section{Definición de integral triple}
%%%%%%%%%%%%%%%%%%%%%%%%%%%%%%%%%%%%%%

\begin{defi}
    Dado un bloque $B$, una partición de orden $n$ es un conjunto
    \[
        P=\{(x_i,y_i,z_i)\in\R^3:i=0,\ldots,n\},
    \]
    donde
    \[
        a=x_0<x_1<\cdots<x_n=b,
    \qquad
        c=y_0<y_1<\cdots<y_n=d
    \]
    y
    \[
        p=z_0<z_1<\cdots<z_n=q.
    \]
    \begin{itemize}
    \item
        A los conjuntos
        \[
            B_{ijk}=[x_{i-1},x_i]\times[y_{j-1},y_j]\times[z_{k-1},z_k],
        \]
        con $i,j,k=1,\ldots,n$, se los llama subbloques de la partición.
    \item
        A las cantidades
        \[
            \Delta x_i=x_i-x_{i-1},
        \qquad
            \Delta y_j=y_j-y_{j-1}
        \texty
            \Delta z_k=z_k-z_{k-1},
        \]
        con $i,j,k=1,\ldots,n$, se las llama el ancho, el alto y la profundidad del subbloque, respectivamente.
    \item
        A la cantidad
        \[
            \Delta V_{ijk} = \Delta x_i\Delta y_j\Delta z_k,
        \]
        con $i,j,k=1,\ldots,n$, se la llama volumen del subbloque.
    \item
        A un conjunto
        \[
            C = \{c^{ijk}\in\R^3: c^{ijk}\in B_{ijk},\ i,j,k=1,\ldots,n\},
        \]
        se lo llama conjunto de etiquetas para $P$.
    \item
        Al par $(P,C)$ se lo llama una partición etiquetada de $B$.
    \item
        A la cantidad
        \[
            |P| = \max_{i=1,\ldots,n}\{\Delta x_i,\Delta y_j,\Delta z_k\}
        \]
        se la llama grosor de la partición.
    \end{itemize}
\end{defi}

\begin{defi}
    Sean $\func{f}{B}{\R}$ un campo escalar y $(P,C)$ una partición etiquetada de $B$, entonces
    \[
        S(f,P,C) = \sum_{i,j,k=1}^n f(c^{ijk})\Delta V_{ijk}
    \]
    es la suma de Riemann de $f$ respecto a la partición etiquetada $(P,C)$.
\end{defi}

\begin{defi}
    Sea $\func{f}{B}{\R}$ un campo escalar, se define la integral triple de $f$ por
    \[
        \iiint_B f\, dV
        =\iiint_B f(x,y,z)\,dx\,dy\,dz
        = \lim_{|P|\to 0} S(f,P,C)
        = \lim_{|P|\to 0} \sum_{i,j,k=1}^n f(c^{ijk})\Delta V_{ijk},
    \]
    siempre que el límite exista. Si el límite existe, se dice que la función es integrable.
\end{defi}

\begin{prop}
    Sea $\func{f}{B}{\R}$ un campo escalar; si $f$ es continua, entonces es integrable.
\end{prop}

\begin{prop}
    Sea $\func{f}{B}{\R}$ un campo escalar; si $f$ es acotada y si el conjunto de discontinuidades de $f$ tiene volumen cero, entonces $\iiint_B f\, dV$ existe.
\end{prop}

%%%%%%%%%%%%%%%%%%%%%%%%%%%%%%%%%%%%%%
\section{Teorema de Fubini para integrales triples}
%%%%%%%%%%%%%%%%%%%%%%%%%%%%%%%%%%%%%%

\begin{teo}
    Sea $\func{f}{B}{\R}$ un campo escalar. Supongamos que $f$ es acotada y que $S$, el conjunto de discontinuidades de $f$, tiene volumen cero. Si cada recta paralela a los ejes de coordenadas interseca a $S$ en una cantidad finita de puntos, entonces
    \begin{align*}
        \iiint_B f\, dV
        & = \int_a^b \l(\int_c^d \l(\int_p^q f(x,y,z)\, dz\r)dy\r)dx \\
        & = \int_a^b \l(\int_p^q\l( \int_c^d f(x,y,z)\, dy\r)dz\r)dx \\
        & = \int_c^d \l(\int_a^b\l( \int_p^q f(x,y,z)\, dz\r)dx\r)dy \\
        & = \int_c^d \l(\int_p^q\l( \int_a^b f(x,y,z)\, dx\r)dz\r)dy \\
        & = \int_p^q \l(\int_a^b\l( \int_c^d f(x,y,z)\, dy\r)dx\r)dz \\
        & = \int_p^q \l(\int_c^d \l(\int_a^b f(x,y,z)\, dx\r)dy\r)dz.
    \end{align*}
\end{teo}

%%%%%%%%%%%%%%%%%%%%%%%%%%%%%%%%%%%%%%
\section{Integrales triples sobre regiones elementales}
%%%%%%%%%%%%%%%%%%%%%%%%%%%%%%%%%%%%%%

\begin{defi}
    Se dice que $W$ es una región elemental en el espacio si puede ser descrita como un subconjunto de $\R^3$ de alguna de las siguientes maneras:
    \begin{itemize}
        \item \textbf{Región de tipo I:} si
        \begin{enumerate}[label=(\alph*)]
            \item $W=\{(x,y,z)\in\R^3 \, : \, \varphi(x,y) \leq z \leq \psi(x,y), \, \gamma(x)\leq y \leq \delta(x),\, a\leq x\leq b \}$, o
            \item $W=\{(x,y,z)\in\R^3 \, : \, \varphi(x,y) \leq z \leq \psi(x,y), \, \alpha(y)\leq x \leq \beta(y),\, c\leq y\leq d \}$.
        \end{enumerate}

        \item \textbf{Región de tipo II:} si
        \begin{enumerate}[label=(\alph*)]
            \item $W=\{(x,y,z)\in\R^3 \, : \, \alpha(y,z) \leq x \leq \beta(y,z), \, \gamma(z)\leq y \leq \delta(z),\, p\leq z\leq q \}$, o
            \item $W=\{(x,y,z)\in\R^3 \, : \, \alpha(y,z) \leq x \leq \beta(y,z), \, \varphi(y)\leq z \leq \psi(y),\, c\leq y\leq d \}$.
        \end{enumerate}

        \item \textbf{Región de tipo III:} si
        \begin{enumerate}[label=(\alph*)]
            \item $W=\{(x,y,z)\in\R^3 \, : \, \gamma(x,z) \leq y \leq \delta(x,z), \, \alpha(z)\leq x \leq \beta(z),\, p\leq z\leq q \}$, o
            \item $W=\{(x,y,z)\in\R^3 \, : \, \gamma(x,z) \leq y \leq \delta(x,z), \, \varphi(x)\leq z \leq \psi(x),\, a\leq x\leq b \}$.
        \end{enumerate}

        \item \textbf{Región de tipo IV:} si $W$ es de tipo I, II y III.
    \end{itemize}
\end{defi}

A partir de aquí, tomaremos a $\Omega$ como un subconjunto cerrado de $B$.

\begin{defi}
    Sea $\func{f}{\Omega}{\R}$ un campo escalar. Definimos la extensión de $f$ a $B$ por
    \[
        f^{\text{ext}}(x,y,z) = \begin{cases}
                          f(x,y,z) & \text{ si } (x,y,z)\in \Omega, \\
                          0 & \text{ si } (x,y,z) \not\in \Omega,
                         \end{cases}
    \]
    para $(x,y,z)\in B$.
\end{defi}

\begin{defi}
    Sea $\func{f}{\Omega}{\R}$ un campo escalar, donde $\Omega$ es una región elemental; si $B$ es cualquier bloque que contiene a $\Omega$, definimos
    \[
        \iiint_\Omega f\, dV = \iiint_B f^{\text{ext}}\, dV.
    \]
\end{defi}

\begin{prop}
    Sea $\Omega$ una región elemental en $\R^3$ y $\func{f}{\Omega}{\R}$ una función continua.
    \begin{itemize}
        \item Si $\Omega$ es de tipo I, entonces
        \[
            \iiint_\Omega f\, dV = \int_a^b \l(\int_{\gamma(x)}^{\delta(x)} \l(\int_{\varphi(x,y)}^{\psi(x,y)} f(x,y,z)\, dz\r)dy\r)dx,
        \]
        o
        \[
            \iiint_\Omega f\, dV = \int_c^d \l(\int_{\alpha(y)}^{\beta(y)} \l(\int_{\varphi(x,y)}^{\psi(x,y)} f(x,y,z)\, dz\r)dx\r)dy.
        \]

        \item Si $\Omega$ es de tipo II, entonces
        \[
            \iiint_\Omega f\, dV = \int_p^q \l(\int_{\gamma(z)}^{\delta(z)} \l(\int_{\alpha(y,z)}^{\beta(y,z)} f(x,y,z)\, dx\r)dy\r)dz,
        \]
        o
        \[
            \iiint_\Omega f\, dV = \int_c^d \l(\int_{\varphi(y)}^{\psi(y)} \l(\int_{\alpha(y,z)}^{\beta(y,z)} f(x,y,z)\, dx\r)dz\r)dy.
        \]

        \item Si $\Omega$ es de tipo III, entonces
        \[
            \iiint_\Omega f\, dV = \int_p^q \l(\int_{\alpha(z)}^{\beta(z)} \l(\int_{\gamma(x,z)}^{\delta(x,z)} f(x,y,z)\, dy\r)dx\r)dz,
        \]
        o
        \[
            \iiint_\Omega f\, dV = \int_a^b \l(\int_{\varphi(x)}^{\psi(x)} \l(\int_{\gamma(x,z)}^{\delta(x,z)} f(x,y,z)\, dy\r)dz\r)dx.
        \]
    \end{itemize}
\end{prop}

%%%%%%%%%%%%%%%%%%%%%%%%%%%%%%%%%%%%%%
\section{Aplicaciones}
%%%%%%%%%%%%%%%%%%%%%%%%%%%%%%%%%%%%%%

\begin{advertencia}
    Dada una región elemental $\Omega$ de $\R^3$, se tiene que $\iiint_\Omega 1\, dV$ es el volumen de la región $\Omega$.
\end{advertencia}

\begin{advertencia}
    Dados una región elemental $\Omega$ de $\R^3$ y $\func{\delta}{\Omega}{\R}$ un campo escalar que representa la densidad en cada punto, el centro de masas $(\overline{x},\overline{y},\overline{z})$ está dado por
    \[
        \overline{x}=
        \frac{\iiint_\Omega x\,\delta(x,y,z)\, dV}
        {\iiint_\Omega \delta(x,y,z)\, dV},
        \quad
        \overline{y}=
        \frac{\iiint_\Omega y\,\delta(x,y,z)\, dV}
        {\iiint_\Omega \delta(x,y,z)\, dV},
        \quad
        \overline{z}=
        \frac{\iiint_\Omega z\,\delta(x,y,z)\, dV}
        {\iiint_\Omega \delta(x,y,z)\, dV}.
    \]
\end{advertencia}

\begin{advertencia}
    Bajo las mismas condiciones, los momentos de inercia con respecto a los ejes coordenados son:
    \[
        I_x = \iiint_\Omega (y^2+z^2)\,\delta\, dV,
        \qquad
        I_y = \iiint_\Omega (x^2+z^2)\,\delta\, dV,
        \qquad
        I_z = \iiint_\Omega (x^2+y^2)\,\delta\, dV.
    \]
\end{advertencia}


\end{document}
